\chapter{Concluzii}
\label{cap:cap5}

Prezenta lucrare a pornit de la o problemă concretă a infrastructurilor de tip cloud privat: deși platforme precum OpenStack oferă un control complet asupra resurselor de calcul, rețea și stocare, complexitatea lor inerentă le face greu accesibile utilizatorului final, pentru care lansarea unei simple mașini virtuale presupune o curbă de învățare abruptă. Obiectivul general, formulat în Capitolul~\ref{cap:cap1}, a constat în proiectarea și implementarea unei platforme intermediare care să abstractizeze această complexitate, oferind o experiență de utilizare simplă, comparabilă cu cea a soluțiilor comerciale, peste o infrastructură OpenStack proprie. În urma dezvoltării și a validării descrise în capitolele anterioare, se poate aprecia că acest obiectiv, împreună cu obiectivele specifice asociate, a fost atins.

Contribuția centrală a lucrării nu constă într-o funcționalitate izolată, ci în \textbf{arhitectura decuplată și orientată pe evenimente} care guvernează întregul sistem. Principiul fundamental conform căruia nivelul API operează exclusiv asupra propriilor baze de date, fără a apela vreodată în mod direct și sincron infrastructura OpenStack, întreaga interacțiune cu aceasta fiind delegată unor procese de fundal prin intermediul evenimentelor de domeniu, a fost decizia de proiectare cu cel mai mare impact. Aceasta asigură simultan receptivitatea interfeței (cererile sunt confirmate imediat, independent de durata reală a provizionării), reziliența sistemului (erorile de infrastructură sunt reconciliate ulterior, nu propagate sincron către utilizator) și testabilitatea, întrucât suprafața API poate fi verificată determinist, separat de infrastructură. Pe acest principiu a fost construită abstractizarea completă a resurselor (rețele, volume, imagini, adrese IP publice), automatizarea întregului ciclu de viață al instanțelor, modelul multi-tenant cu control al accesului bazat pe roluri și cu cote configurabile, precum și consola de administrare la nivel de platformă.

În raport cu soluțiile existente analizate în Capitolul~\ref{cap:cap2}, platforma dezvoltată ocupă o poziție intermediară pe care niciuna dintre acestea nu o acoperă pe deplin. Spre deosebire de interfața nativă OpenStack (Horizon), proiectată pentru operatori de sistem și încărcată de terminologie de infrastructură, soluția propusă ascunde această complexitate în spatele unor fluxuri de lucru simple. Spre deosebire de platformele comerciale închise, precum panourile de tip Fleio, ea este o soluție open source și adaptabilă, potrivită cloud-urilor private și mediilor educaționale. Deși se inspiră din experiența de utilizare a platformelor precum DigitalOcean, platforma dezvoltată oferă această simplitate peste o infrastructură controlată local, fără dependența de un furnizor public.

Validarea funcțională, prezentată în Capitolul~\ref{cap:cap4}, a confirmat corectitudinea soluției pe un cluster OpenStack real, desfășurat prin Kolla-Ansible: provizionarea instanțelor se finalizează cu succes, cheile SSH generate de platformă sunt injectate corect și permit accesul efectiv la mașini, iar consola noVNC oferă un canal de administrare alternativ, independent de conectivitatea la nivel de rețea. Aceeași testare a evidențiat limitările soluției, ținând în principal de mediul de desfășurare pe un singur nod, fără redundanță, și de izolarea realizată la nivel logic, nu la nivel de infrastructură.

\section{Direcții viitoare de dezvoltare}
\label{cap:cap5:directii-viitoare}

Arhitectura decuplată a platformei lasă deschise mai multe direcții de extindere, care depășesc cadrul prezentei lucrări fără a-i contrazice premisele.

O primă direcție vizează \textbf{disponibilitatea înaltă}: desfășurarea OpenStack pe mai multe noduri, într-o configurație redundantă, ar elimina punctul unic de defecțiune al topologiei \textit{all-in-one}, transformând platforma dintr-o soluție de dezvoltare într-una aptă pentru producție. O \textbf{izolare la nivel de infrastructură} ar putea fi obținută prin alocarea câte unui proiect OpenStack distinct fiecărei organizații, înlocuind izolarea logică actuală cu una garantată de Keystone.

O a doua direcție privește îmbogățirea funcționalităților oferite utilizatorului: \textbf{încărcarea directă a imaginilor personalizate} (prin transmiterea fișierului prin platformă către Glance, alături de importul de la un URL deja implementat), \textbf{redimensionarea instanțelor} (modificarea pachetului hardware al unei mașini existente) și îmbunătățirea sistemului de \textbf{telemetrie și metering} către un model complet de facturare. Deoarece procesarea sarcinilor este deja delegată unui broker de mesaje, \textbf{scalarea orizontală} a procesării asincrone prin rularea mai multor workeri Celery concurenți pe un sistem Linux ar fi imediată, fără modificări de arhitectură.

\section{Lecții învățate pe parcursul dezvoltării proiectului de diplomă}
\label{cap:cap5:lectii}

Dezvoltarea proiectului a confirmat, în mod practic, câteva principii care la început păreau mai degrabă teoretice. Cel mai important este că o \textbf{decuplare riguroasă a responsabilităților} nu reprezintă doar o eleganță de proiectare, ci aduce beneficii concrete și măsurabile: aceeași regulă care interzice apelurile directe către OpenStack din nivelul API a făcut posibilă, ulterior, testarea automată a întregului flux funcțional, întrucât logica de business a putut fi verificată independent de infrastructură.

De asemenea, integrarea cu o platformă complexă precum OpenStack a evidențiat diferența dintre corectitudinea aplicației și starea reală a infrastructurii subiacente: erori precum lipsa unui nod disponibil pentru planificare (\texttt{NoValidHost}) sau indisponibilitatea unui serviciu după suspendarea mașinii gazdă nu țin de codul aplicației, dar trebuie anticipate și tratate. Tocmai de aceea, mecanismele de reconciliere periodică a stărilor și de înregistrare a evenimentelor s-au dovedit esențiale pentru a menține o stare consistentă a sistemului. Gestionarea evoluției schemei bazei de date, în absența unui sistem dedicat de migrații, a subliniat dificultatea modificărilor de structură pe un sistem aflat deja în funcțiune și importanța unei proiectări foarte minuțioase a aplicației înainte de implementarea acesteia.
